% Figure from MATH 493 notes, section 8. Compile with the notes preamble.
\begin{tikzpicture}[
    pt/.style={circle, draw, thick, minimum size=14pt, inner sep=0pt, fill=white, font=\scriptsize},
    arr/.style={-{Stealth[length=4pt]}, thin, shorten >=1pt, shorten <=1pt}]
  % left: x -> 3x on the nonzero classes mod 7
  \begin{scope}
    \foreach \x in {1,...,6} {
      \node[pt, draw=figB] (L\x) at (0,{-0.55*\x}) {$\x$};
      \node[pt, draw=figB] (R\x) at (2.2,{-0.55*\x}) {$\x$};
    }
    \foreach \x/\y in {1/3, 2/6, 3/2, 4/5, 6/4} \draw[arr, figB] (L\x) -- (R\y);
    \draw[arr, figR, thick] (L5) -- (R1);
    \node[pt, draw=figR, fill=figR!15] at (R1) {$1$};
    \node[font=\small, figB] at (1.1,0.05) {$x \mapsto 3x \bmod 7$};
    \node[font=\scriptsize, figR, anchor=west] at (2.45,-0.55) {$3 \cdot 5 \equiv 1$};
  \end{scope}
  % right: x -> 4x on Z/6
  \begin{scope}[xshift=4.9cm]
    \foreach \x in {0,...,5} {
      \node[pt, draw=black!60] (L\x) at (0,{-0.55*(\x+1)}) {$\x$};
    }
    \foreach \x in {0,2,4} \node[pt, draw=black!60] (R\x) at (2.2,{-0.55*(\x+1)}) {$\x$};
    \foreach \x in {1,3,5} \node[pt, draw=figR, dashed, text=figR, thin] (R\x) at (2.2,{-0.55*(\x+1)}) {$\x$};
    \foreach \x/\y in {0/0, 1/4, 2/2, 3/0, 4/4, 5/2} \draw[arr, black!60] (L\x) -- (R\y);
    \node[font=\small] at (1.1,0.05) {$x \mapsto 4x \bmod 6$};
    \node[font=\scriptsize, figR, anchor=west] at (2.45,-1.1) {never hit};
  \end{scope}
\end{tikzpicture}
